\section{Stationary boundary profiles and their constants}\label{st:profile-section}
The derivative of a fixed-output profile describes how the
maximum changes when the output moves. It is convenient to
record this derivative together with the output itself.
For a strict open output $b$, define
\begin{equation}\label{st:port}
 B=\cot\pi b,\qquad \kappa_E=e^{-f_E'(b)},\qquad
 X=\sqrt{\frac{1+B^2}{1+\kappa_E}},\qquad
 Y=\sqrt{\kappa_E}\,X .
\end{equation}
The square roots are positive. Thus
$X^2+Y^2-B^2=1$ and $f_E'(b)=2\log(X/Y)$.
We call $(B,X,Y)$ the stationary boundary parameters.
Their definitions are variational; no global geometric
realization is assumed.

\begin{lemma}[The stationary product]\label{st:port-product}
At a strict stationary product the two children and the
output have the same $X$. With child cotangents $B_1,B_2$,
\begin{equation}\label{st:port-product-equations}
 S=B_1+B_2>0,\qquad
 B_o=\frac{B_1B_2+X^2-1}{S},\qquad
 Y_o=\frac{Y_1Y_2}{S}.
\end{equation}
A turn sends $(B,X,Y)$ to $(-B,Y,X)$.
\end{lemma}
\begin{proof}
At inputs $x,y$ and output $z$, stationarity gives
$f_1'(x)=-m_x$, $f_2'(y)=-m_y$ and $f_o'(z)=m_z$.
Differentiating \eqref{st:m} yields
\begin{align*}
 e^{m_x}&=\frac{\sin\pi(x+y)\sin\pi(z-x)}
                  {\sin\pi x\sin\pi(x+y-z)},\\
 e^{m_y}&=\frac{\sin\pi(x+y)\sin\pi(z-y)}
                  {\sin\pi y\sin\pi(x+y-z)},\\
 e^{-m_z}&=\frac{\sin\pi(z-x)\sin\pi(z-y)}
                  {\sin\pi z\sin\pi(x+y-z)}.
\end{align*}
Set $D=1-B_1B_2+B_o(B_1+B_2)$. The identities
\[
 S=\frac{\sin\pi(x+y)}{\sin\pi x\sin\pi y}>0,\qquad
 D=\frac{\sin\pi(x+y-z)}
          {\sin\pi x\sin\pi y\sin\pi z}>0
\]
and \eqref{st:port} give
\[
 X_1^2=X_2^2=X_o^2=D,\quad
 Y_1^2=S(B_1-B_o),\quad Y_2^2=S(B_2-B_o),\quad
 Y_o^2=(B_1-B_o)(B_2-B_o).
\]
Taking positive roots proves the product formulas.
For a turn $f_{\T E}(b)=f_E(1-b)$, so its derivative
changes sign and its cotangent changes sign. This
proves the last assertion.
\end{proof}

\begin{proposition}[Parameters of the three types]\label{st:port-types}
On their actual strict output intervals the parameters satisfy
\begin{equation}\label{st:parameter-table}
 \begin{array}{c|c}
 L&(B,X,Y)=(B,B,1),\quad B>0,\\
 H&(B,X,Y)=(-Y,1,Y),\quad Y>0,\\
 A&-Y<B<X.
 \end{array}
\end{equation}
Consequently there are constants $c_E$ such that
\begin{equation}\label{st:rigid-profiles}
 f_E(b)=g(b)+c_E\quad(E:L),\qquad
 f_E(b)=g(1-b)+c_E\quad(E:H).
\end{equation}
\end{proposition}
\begin{proof}
For a leaf, $g'=2\log\cot\pi b$ gives $(B,B,1)$.
The turn formula exchanges the first two rows and
preserves the inequality in the third.
Consider the products giving type $A$, whose children
have type $L$ or $A$.
If the common $X<1$, put $h_0=\sqrt{1-X^2}$ and write
\[
 B_i=h_0\cot\alpha_i,\qquad
 Y_i=h_0\csc\alpha_i,\qquad 0<\alpha_i<\pi.
\]
The low type has $\alpha_i=\arccos X$ and the
unrestricted type has $\alpha_i>\arccos X$.
The positivity of $S$ gives $\alpha_1+\alpha_2<\pi$,
and \eqref{st:port-product-equations} adds these angles.
Therefore $B_o<X$ and
$B_o+Y_o=h_0\cot((\alpha_1+\alpha_2)/2)>0$.
If $X\ge1$, each child has
$\sqrt{X^2-1}<B_i\le X$, with equality at the
upper bound allowed for type $L$.
In particular $B_o>0$, and
\[
 (X-B_1)(X-B_2)<1
\]
implies $B_o<X$ by \eqref{st:port-product-equations}.
The strict inequality holds also at $X=1$, since
each $B_i>0$.

A product giving type $H$ has one high child, with
parameters $(-u,1,u)$. Hence the common $X=1$.
Its other child, of type $L$ or $A$, has parameters
$(b,1,b)$ with $b>0$. Strictness says $b-u>0$,
and the output is
\[
 \left(-\frac{ub}{b-u},1,\frac{ub}{b-u}\right).
\]
This proves all cases of the table.
Finally the first two rows determine $f_E'$ as
the derivative of $g(b)$ or $g(1-b)$ on the
connected strict interval, giving \eqref{st:rigid-profiles}.
Their constants are computed next, without a choice
of additive normalization.
\end{proof}

\subsection{The intrinsic satellite action}
No product has type $L$. A product of type $H$ has
exactly one high child. Starting at an $L$ or $H$
output, follow turns and the unique high child of
each product. This path terminates at a unique
original leaf, the \emph{distinguished leaf}.
The other children on the path are called its
satellites, and have type $L$ or $A$.
Between successive high products the number of turns
is even. Thus the full algebra also gives
\begin{equation}\label{st:spine-source}
 E_H=(\T h_*)\prod_{j=1}^m Q_j,\qquad
 E_L=\T\left((\T h_*)\prod_{j=1}^m Q_j\right),
\end{equation}
after associating the inherited products. This
identity retains the actual polynomials; its
stationary profile is simpler than the polynomial.

For a satellite $Q$ define
\begin{equation}\label{st:Omega}
 \Omega_Q=\max_{0\le a\le1/2}\{g(a)+f_Q(a)\}.
\end{equation}
The maximum is the complete action of the literal
strict closure $C(h\,\T Q)$. The strict type of
$Q$ is $L$ or $A$, so Theorem~\ref{st:positive-action}
gives a unique strict maximizer $\beta_Q$.
Its equation is
$g'(\beta_Q)+f_Q'(\beta_Q)=0$.
Consequently the satellite has parameters
\begin{equation}\label{st:satellite-parameters}
 (B_Q,X_Q,Y_Q)=(b_Q,1,b_Q),\qquad
 b_Q=\cot\pi\beta_Q>0 .
\end{equation}
For type $L$, $\beta_Q=1/4$ and $b_Q=1$;
for type $A$, the strict inequality $B_Q<X_Q$
gives $1/4<\beta_Q<1/2$ and $0<b_Q<1$.

\begin{theorem}[Harmonic response and absolute profiles]\label{st:profiles}
For the spine \eqref{st:spine-source}, let $u_0$
be the cotangent of its distinguished original leaf.
Let $v$ be the final high output's $Y$, equivalently
the final low output's $X$. At its stationary state,
\begin{equation}\label{st:harmonic}
 \frac1{u_0}=\frac1v+\sum_{j=1}^m\frac1{b_{Q_j}},
 \qquad
 c_E=\sum_{j=1}^m\Omega_{Q_j}.
\end{equation}
In particular \eqref{st:rigid-profiles} holds with
exactly this constant. The parameters $b_{Q_j}$
depend only on the corresponding satellites.
\end{theorem}
\begin{proof}
Every high product has common $X=1$, so its satellite
is at the unique state \eqref{st:satellite-parameters}.
With high input $(-u,1,u)$ and output $(-v,1,v)$,
\eqref{st:port-product-equations} gives
\begin{equation}\label{st:one-response}
 v=\frac{ub}{b-u},\qquad
 0<u<b,\qquad \frac1u=\frac1v+\frac1b .
\end{equation}
Iteration proves the first formula. Conversely every
$v,b>0$ gives the strict branch $0<u<b$; no sign
has been lost by squaring the stationary equations.

It remains to fix the constant. Suppose the incoming
high profile is $g(1-x)+c$. Send the outgoing high
label $z$ to one. Then
$v=-\cot\pi z\to+\infty$, $u\to b$ and
$x\to1-\beta_Q$. The last gate approaches
$x+\beta_Q=z=1$. Its continuous action is zero:
in \eqref{st:m} use $J(1-\beta_Q)=-J(\beta_Q)$,
or note that all three atoms in
\eqref{st:gate-atoms} are flat.
Continuity of profiles, proved in
Theorem~\ref{st:positive-action}, gives
\[
 c_{\rm out}=c+g(\beta_Q)+f_Q(\beta_Q)=c+\Omega_Q.
\]
Turns preserve the constant and the bare leaf has
constant zero. This proves the second formula.
\end{proof}
For a low satellite, $\Omega_Q=2G_0+c_Q$ where
$G_0=g(1/4)$. The formula is an identity of optimized
actions. For example at $N=3$ the complete source has
$h=10-z$ and $\T h=z^2-1$; their product is
$10z^2-2z-8$ in
$\CC[z]/((z-1)^2(z+2))$. It is not a scalar
multiple of $\T h$: the value ratio at $-2$ is
$12$, whereas the normal ratio at $1$ is $9$.
Thus the constant profile does not suppress normal
data of the source.

\subsection{Unrestricted components and their effective range}
Replace each maximal low subtree by its profile
$g+c_i$. A remaining type-$A$ subtree has syntax
\begin{equation}\label{st:skeleton}
 A::=LL\mid\T A\mid AA\mid LA\mid AL,
\end{equation}
where $L$ here is an effective profile, not a
replacement quantum leaf. The definition of the
maximum immediately gives
\begin{equation}\label{st:constant-skeleton}
 f_A(z)=f_{A^{\rm skel}}(z)+\sum_i c_i.
\end{equation}
At an unflagged junction directed toward a parent,
the outgoing expression is $Q=\T A$ with $A=FG$
an unturned product and both children of type $L$
or $A$. Its intrinsic action is
\[
 \Omega_Q=\max_{0\le a\le1/2}
                 \{g(a)+f_A(1-a)\}.
\]
This is the action of the strict closure $C(hA)$.

\begin{theorem}[The intrinsic probe range]\label{st:probe}
At this closing maximum the unturned output has
\begin{equation}\label{st:probe-port}
 (B_A,X_A,Y_A)=(-x,x,1),\qquad \frac12\le x<1.
\end{equation}
Every effective leaf in \eqref{st:skeleton},
including the closing probe, has cotangent in
$[1/2,1)$. Equality at $1/2$ occurs only when
the entire effective closure consists of three
leaves. In every other case all these cotangents
are strictly between $1/2$ and one.

The absolute identities are
\begin{equation}\label{st:probe-constant}
 \Omega_Q=F_{\rm probe}(A^{\rm skel})+\sum_i c_i,
 \qquad F_0=\sum_j\Omega_{Q_j}
\end{equation}
for a whole reduced tree rooted at a flagged leaf,
with satellites $Q_j$ at that root.
\end{theorem}
\begin{proof}
The root condition says that the two output labels
are complementary and their profile derivatives
agree. Hence their triples are $(B,X,Y)$ and
$(-B,X,Y)$. The bare probe has $(x,x,1)$,
so \eqref{st:probe-port} holds initially with
$0<x<1$, the strict upper bound coming from
$-Y_A<B_A$.
At the last product use the identities in the
proof of Lemma~\ref{st:port-product}:
\[
 1=Y_A^2=(B_F-B_A)(B_G-B_A)
             \le(X_A-B_A)^2=4x^2 .
\]
Both factors are positive because output labels
strictly exceed the inputs. This proves $x\ge1/2$.

At the root product every child has common $X=x<1$.
Its cotangent satisfies $B_i\le X$ and, by
$B_1+B_2>0$, also $B_i>-X$.
If it is type $A$, both inequalities are strict,
so $Y_i^2=1+B_i^2-X^2<1$.
Turns exchange $X,Y$. The same argument propagates
the strict bounds $X,Y<1$ through every proper
type-$A$ descendant.
For a low child $(X,X,1)$ and another child
$(B,X,Y)$, the output bound $Y_o\le1$ implies
\begin{equation}\label{st:effective-bound}
 1+B^2-X^2\le(B+X)^2,\qquad
 1\le2X(B+X)\le4X^2.
\end{equation}
The last inequality is strict for an unrestricted
neighbor. For two low children at a proper
descendant, the first is strict because $Y_o<1$.
Thus equality is possible only at the root
product with two effective low children. It
gives the three-leaf closure and $x=1/2$ at
all three leaves. This proves the range.

The first identity in \eqref{st:probe-constant}
is \eqref{st:constant-skeleton} inside the same
maximum. For the second, write the original
root source as $\phi((\T h_0)\prod Q_j)$.
The associated open high profile is
$g(1-z)+\sum_j\Omega_{Q_j}$ by
Theorem~\ref{st:profiles}. At the closing output
$z=1$ the deleted last gate has coefficient
one and $g(0)=0$. Profile continuity identifies
its endpoint maximum with precisely the original
closed action, proving the formula.
\end{proof}

One boundary case in later coefficient estimates can
be read directly from these parameters. If the last
product is $A=LF$ with $F$ unrestricted, then
\[
 B_F=\frac1{2x}-x,\qquad Y_F=\frac1{2x},\qquad
 f_F'=2\log(2x^2),\quad f_A'=2\log x .
\]
For $1/2\le x\le2/3$ both slopes lie in
$[-2\log2,0]$. This is an exact closed parameter
range; it does not assert an estimate for the
signed source. Such estimates require its actual
coefficient arrays.

\begin{corollary}[The action budget at any exposed leaf]\label{st:root-budget}
Use Corollary~\ref{st:any-root}, with root edge parity
$\sigma$ and remaining product $A$.
If $\sigma=1$, every branch of $A$ has type $L$ or
$A$ and the original complete action is
$F_0=\sum_Q\Omega_Q$, summed over those branches.
If $\sigma=0$, then $Q=\T A$ has type $L$ or $A$
and $\Omega_Q=F_0$.
\end{corollary}
\begin{proof}
For $\sigma=1$, this is the last identity of
\eqref{st:probe-constant} at the exposed flagged
root. For $\sigma=0$, the original positive closure
is $C(hA)$, whose action is exactly
$\max_a\{g(a)+f_A(1-a)\}=\Omega_{\T A}$.
This includes the decorated case in which $A$
has one high child and $\T A$ is low.
Only the finite source is rerooted; the variational
maximum and every physical leaf are unchanged.
\end{proof}

\section{The quarter bound at the original closed maximum}
\begin{theorem}[All original leaves are at or above quarter]\label{st:quarter}
At the unique complete maximum of any strict closed
expression, every stationary port satisfies $X,Y\le1$.
Every original leaf label therefore satisfies
\begin{equation}\label{st:quarter-bound}
 \frac14\le a_\ell(y_*)<\frac12 .
\end{equation}
Along an $L/H$ spine with $m$ satellites, its
distinguished original leaf additionally satisfies
$\cot(\pi a_*)\le1/(m+1)$.
\end{theorem}
\begin{proof}
At the complementary root the two triples have
opposite $B$ and equal $X,Y$. For type $AA$ the
inequalities give $|B|<\min(X,Y)$.
Together with $X^2+Y^2-B^2=1$ this gives
$X,Y<1$. For type $AL$, the low side is
$(x,x,1)$ and the unrestricted inequality
$-1<-x$ gives $x<1$.
For type $AH$, the high side has $X=1$,
and the unrestricted side gives $Y<1$.
For type $LH$, the two explicit forms force
$X=Y=1$. These exhaust the strict closing pairs.

Descend through a product. The common $X$ is
already at most one. Every child type satisfies
$B_i\le X$, and $B_1+B_2>0$ implies
$B_i>-B_j\ge-X$. Thus
$Y_i^2=1+B_i^2-X^2\le1$.
A turn exchanges the same two bounds.
At a leaf, $X=\cot\pi a>0$, proving
\eqref{st:quarter-bound}. Finally, in
\eqref{st:harmonic} one has $v\le1$ and
$b_{Q_j}\le1$, so $1/u_0\ge m+1$.
\end{proof}
The lower endpoint is attained. In
$C(h\,\T(h^s))$, the inner $s$ original leaves
have $X=1$ and hence label $1/4$.
Indeed their product has at $X=1$ the parameters
$B=Y=1/s$, its turn has $(-1/s,1/s,1)$, and
the closing leaf has cotangent $1/s$.
Thus strictness of every angle state does not force
its original leaf labels to lie strictly above one quarter.
